\documentclass[aps,prx,reprint,superscriptaddress,nofootinbib]{revtex4-2}

\usepackage{amsmath,amssymb,amsthm}
\usepackage{physics}
\usepackage{graphicx}
\usepackage{hyperref}
\usepackage{xcolor}
\usepackage{enumitem}

\hypersetup{
    colorlinks=true,
    linkcolor=blue,
    citecolor=teal,
    urlcolor=blue
}

% Theorem Environments
\newtheorem{theorem}{Theorem}
\newtheorem{lemma}[theorem]{Lemma}
\newtheorem{corollary}[theorem]{Corollary}
\theoremstyle{definition}
\newtheorem{definition}{Definition}
\newtheorem{observation}{Observation}

\begin{document}

\title{Topological Determinism and the Dynamical Shift-Rank Bound:\\ A Multi-Level Compiler Architecture for Universal Quantum Simulation}

\author{THIS IS AI GENERATED FOR NOW}
\email{you@unitary.foundation}
\affiliation{Unitary Foundation}
% Add other authors as needed

\date{\today}

\begin{abstract}
The exact simulation of universal fault-tolerant quantum circuits is fundamentally bottlenecked by the exponential complexity of non-Clifford operations. Recent advances utilizing the generalized stabilizer formalism have extended classical simulation limits by tracking sparse superpositions over dynamic stabilizer bases. However, these methods rely on static spatial bounds and perform $\mathcal{O}(n^2)$ matrix updates at runtime, severely constraining performance. In this work, we introduce the Universal Compiler Collection (UCC), a framework that fundamentally decouples the Boolean geometry of a quantum state from its probabilistic execution. By formalizing a property we term \textit{Topological Determinism}, we shift all topological matrix tracking into an Ahead-Of-Time (AOT) compiler. Operating in this decoupled framework, we prove the \textit{Dynamical Shift-Rank Theorem}, demonstrating that the active support of the continuous state vector forms an exact affine subspace over $\mathbb{F}_2^n$. This establishes a strictly tighter, dynamically evolving bound on the memory footprint of universal circuits, naturally capturing temporal collisions and active measurement cooling that spatial bounds fail to resolve. We interpret the dimension of this affine space as the ``Active Magic Volume'' of the circuit. Furthermore, we outline the UCC architecture—featuring FTQC-safe symplectic optimization and a cache-aligned virtual machine—and discuss its application to ground-truth simulations of distance-5 magic state cultivation protocols at unprecedented scales.
\end{abstract}

\maketitle

% ==============================================================================
\section{Introduction and the Principle of Topological Determinism}
\label{sec:intro}
% ==============================================================================

The transition from noisy intermediate-scale quantum (NISQ) devices to fault-tolerant quantum computing (FTQC) requires simulating circuits with millions of operations. These circuits are heavily dominated by Clifford gates, stabilizer measurements, and a sparse distribution of non-Clifford logic (e.g., $T$-gates). By the Gottesman-Knill theorem \cite{gottesman1998}, pure Clifford circuits can be simulated efficiently. However, injecting universal non-Clifford operations causes the required computational resources to scale exponentially \cite{bravyi2016}.

Recent state-of-the-art approaches, such as the SOFT simulator \cite{li2025soft}, utilize the generalized stabilizer formalism \cite{yoder2012} to simulate these circuits by tracking a stabilizer tableau alongside a sparse array of continuous amplitudes. While theoretically elegant, existing implementations natively couple the geometric tableau evolution with runtime execution. This coupling creates a severe performance bottleneck: evaluating $\mathcal{O}(n^2)$ dense Boolean matrix operations dynamically during Monte Carlo sampling. Furthermore, previous bounds on the required size of the state vector rely on static, spatial constraints that vastly overestimate the complexity of highly entangled circuits.

In this work, we resolve these limitations through a new multi-level compiler architecture based on a physical property we formalize as \textbf{Topological Determinism}.

\begin{observation}[Topological Determinism]
In any quantum circuit composed of Clifford unitaries, non-Clifford superpositions, arbitrary Pauli noise channels, and projective Pauli measurements (lacking classically controlled Clifford gates), the Boolean geometry of the stabilizer frame is strictly independent of stochastic noise events, random measurement outcomes, and continuous non-Clifford amplitudes.
\end{observation}

Because stochastic Pauli noise merely flips scalar $\pm 1$ signs, and non-Clifford gates populate orthogonal branches relative to a shared geometry, the Boolean matrix grid defining the basis never warps randomly. Projective measurements induce Aaronson-Gottesman (AG) row reductions whose Boolean trajectories are deterministically fixed prior to the random outcome \cite{aaronson2004}.

Therefore, the entire sequence of generalized tableaus $\{ \mathcal{T}_t \}$ can be deterministically computed offline by an Ahead-Of-Time (AOT) compiler. The runtime Virtual Machine (VM) need only manage the continuous probability amplitudes $\mathbf{v}_t$ and a classical sign tracker to account for physical noise parity relative to the pre-compiled reference frame. This decoupled framework allows us to rigorously map the evolution of the continuous amplitudes to an affine subspace over a finite field, yielding a novel, strictly tight bound on simulation complexity.
% ==============================================================================
\section{Dual-Picture Decoupling and the Static Affine Subspace}
\label{sec:algebra}
% ==============================================================================

\subsection{The Limitation of Forward-Evolving Tableaus}

In the standard generalized stabilizer formalism \cite{yoder2012, li2025soft}, the quantum state at step $t$ is represented as a sparse superposition over a moving orthonormal basis induced by the current, forward-evolving tableau $\mathcal{T}_t = (\mathcal{S}_t, \mathcal{D}_t)$:
\begin{equation}
    |\phi_t\rangle = \sum_{\alpha \in \Omega_t} v_\alpha |b_{t,\alpha}\rangle \quad \text{where} \quad |b_{t,\alpha}\rangle = d_{t}^\alpha |\psi_{\mathcal{S}_t}\rangle
\end{equation}
While mathematically exact, implementing this directly forces classical simulators into a conceptual ``Schr\"odinger trap.'' By coupling the continuous amplitudes $v_\alpha$ to a forward-evolving geometric frame $|b_{t,\alpha}\rangle$, the runtime execution thread is forced to carry the $\mathcal{O}(n^2)$ Boolean matrix tracking along for the ride. Because a physical circuit contains operations that are fundamentally stochastic (e.g., random measurement outcomes and Pauli noise channels), this approach binds heavy geometric matrix updates directly to the probabilistic Monte Carlo hot-loop, creating a severe computational bottleneck.

\subsection{Topological Determinism and Heisenberg Rewinding}

We break this bottleneck by applying the principle of \textbf{Topological Determinism}. Let a quantum circuit be defined as a sequence of operations $\mathcal{C} = (O_0, O_1, \dots, O_{N-1})$, where each operation $O_t$ may be a Clifford unitary, a non-Clifford rotation, a projective Pauli measurement, or a probabilistic Pauli noise channel.

\begin{observation}[Topological Determinism]
In the absence of classically controlled Clifford gates, the geometric trajectory of the circuit's Boolean reference frame is strictly independent of stochastic noise events, continuous amplitude branches, and random measurement outcomes.
\end{observation}

While the physical state trajectory is highly probabilistic, its underlying Boolean geometry is entirely predetermined. Pauli noise merely flips scalar $\pm 1$ parity tracking signs without altering the Pauli strings. Non-Clifford rotations spawn orthogonal branches relative to a shared coordinate system. Projective measurements physically collapse the state, but the Aaronson-Gottesman (AG) matrix row-reductions \cite{aaronson2004} required to find the new commuting basis are strictly deterministic functions of the Pauli algebra, independent of the random $\pm 1$ measurement outcome.

Because the topological trajectory is completely decoupled from the stochastic execution, we can separate them architecturally: evaluating the geometry Ahead-Of-Time (AOT) and the probability at runtime. Rather than pushing the tableau forward, an AOT compiler can use the Heisenberg picture to pull the operations backward (conceptually mirroring the reverse-tableau tracking of \texttt{Stim} \cite{gidney2021stim}).

To maintain mathematical rigor across non-unitary operations (measurements and noise), we avoid standard unitary conjugation notation. Instead, let $\Phi_{t \to 0}$ denote the cumulative, deterministic geometric transformation from step $t$ back to $t=0$. This symplectic mapping explicitly tracks the inverse sequence of Clifford conjugations and deterministic AG basis changes, while completely ignoring scalar noise flips and continuous amplitudes. For any physical operation $O_t$ acting along a Pauli axis $P_t$, we rewind the operator to the static $t=0$ reference frame $\mathcal{T}_0 = (\mathcal{S}_0, \mathcal{D}_0)$:
\begin{equation}
    P_{t \to 0} = \Phi_{t \to 0}(P_t)
\end{equation}
We then decompose this rewound operator against the initial reference frame:
\begin{equation}
    P_{t \to 0} \propto \prod_{i=1}^n (d_{0,i})^{\beta_i(P)} (s_{0,i})^{\gamma_i(P)}
\end{equation}
where $\beta(P) \in \mathbb{F}_2^n$ is the $t=0$ \textit{spatial shift mask} and $\gamma(P) \in \mathbb{F}_2^n$ is the $t=0$ \textit{phase mask}.

This mathematical rewinding completely absorbs all Clifford operations. The original circuit $\mathcal{C}$ is flattened into a static, geometry-free list of Pauli operations defined strictly by their pre-computed masks. In our architecture, this formulation serves as the formal mathematical definition of the \textbf{Heisenberg Intermediate Representation (HIR)}.

\subsection{Global Commutation and Dual-Mode Compilation}

A profound consequence of flattening the circuit into the $t=0$ basis is the emergence of global commutativity. In a standard forward-evolving simulation, determining whether a Pauli operation $P_{t_1}$ commutes with a later operation $Q_{t_2}$ requires evaluating the entire intervening sequence of Clifford gates. In the HIR, both operators are natively rewound to the exact same static vacuum.

Because the geometric mapping $\Phi_{t \to 0}$ is a symplectic transformation, it strictly preserves the algebraic structure of the Pauli group. The temporal commutation relation of the original operations is exactly equivalent to the static commutation of their rewound masks. This translates deep temporal dependencies into a single $\mathcal{O}(1)$ bitwise symplectic inner product between their respective $(\beta, \gamma)$ masks, evaluated purely in the $t=0$ stabilizer frame.

This algebraic flattening formally motivates the architectural distinction between the HIR and the execution engine, revealing a powerful duality in our framework: \textit{optimization is independent of the state, but execution requires a known state.}

This naturally enables two distinct compilation pathways:
\begin{itemize}[leftmargin=*]
    \item \textbf{State-Agnostic Optimization (Hardware Snippets):} By choosing \textit{not} to rely on (or collapse) a specific $t=0$ stabilizer tableau (e.g., initializing the compiler with a generic Identity tableau), we unlock state-independent compilation. A circuit ``snippet'' or logic subroutine can be traced and optimized aggressively. The compiler's Middle-End can fuse, cancel, and commute operations globally across the snippet purely via their $t=0$ masks, completely blind to the eventual input state. This optimized block can then be directly lowered to target physical hardware (e.g., as a Pauli-Based Computation routing template) or seamlessly composited into larger algorithms.

    \item \textbf{State-Aware Compilation (Simulation):} Conversely, when a full algorithm is executed via Monte Carlo simulation (or when executing on a specific on-device state), we trace the circuit from a known initial state (such as $|0\rangle^{\otimes n}$). The compiler actively tracks deterministic measurement collapses and anchors the spatial shift vectors $\beta$, fully resolving the exact boundary conditions required by the runtime Virtual Machine.
\end{itemize}

\subsection{Execution on the Static Vacuum}

When operating in the State-Aware mode, the geometry is completely flattened and optimized in the HIR, freeing the runtime Virtual Machine (VM) from topological tracking. The VM evaluates the continuous probability amplitudes forward in time (the Schr\"odinger picture), but it does so entirely over the \textit{fixed} $t=0$ state-aware vacuum basis:
\begin{equation}
    |\phi_t\rangle = \sum_{\alpha \in \Omega_t} v_\alpha |b_{0,\alpha}\rangle
\end{equation}
The runtime needs only to manage the continuous array $\mathbf{v}_t$ and a lightweight classical bit-vector tracking the $\pm 1$ physical phase inversions generated by stochastic noise. A rewound Pauli operator from the HIR acts on this static basis as a spatial shift up to a relative phase:
\begin{equation}
    P_{t \to 0} |b_{0,\alpha}\rangle = \xi_\alpha(P) |b_{0,\alpha \oplus \beta(P)}\rangle
\end{equation}
where $\xi_\alpha(P) \propto (-1)^{\gamma(P) \cdot \alpha}$. What were previously $\mathcal{O}(n^2)$ matrix updates in a dynamic tableau are now reduced to $\mathcal{O}(1)$ bitwise operations and fast XOR array indexing.

\subsection{The Dynamical Shift-Rank Bound}

Prior literature \cite{li2025soft} bounded the required sparsity of the coefficient vector $|\mathbf{v}_t|$ by the static spatial support of the non-Clifford gates. However, this static bound is blind to the temporal dynamics of the circuit. By framing the execution over the static $t=0$ vacuum, we reveal a strictly tighter, dynamically exact bound.

\begin{definition}[The Active Support]
Let $\Omega_t \subseteq \mathbb{F}_2^n$ be the support of the coefficient vector over the static $t=0$ basis, defined as the set of binary keys where the amplitude $v_\alpha \neq 0$.
\end{definition}

\begin{theorem}[The Dynamical Shift-Rank Bound]
\label{thm:shift_rank}
At any step $t$ of a universal quantum circuit, the state vector support $\Omega_t$ forms an exact affine subspace over the static vacuum $\mathbb{F}_2^n$. That is, $\Omega_t = \alpha_t \oplus V_t$, where $V_t$ is a linear vector space over $\mathbb{F}_2$ and $\alpha_t$ is a translation vector.

Consequently, the number of non-zero continuous amplitudes is strictly bounded by:
\begin{equation}
    |\mathbf{v}_t| \le 2^{\dim(V_t)}
\end{equation}
where $\dim(V_t)$ evolves dynamically based purely on the $\mathbb{F}_2$ linear dependence of the rewound Pauli shift masks $\beta(P)$, independent of the continuous amplitudes.
\end{theorem}

\begin{proof}
We proceed by structural induction on the operations of the state-aware HIR acting on the static vacuum. At $t=0$, the system is in the vacuum state $|\phi_0\rangle = |b_{0,\vec{0}}\rangle$. The support is $\Omega_0 = \{\vec{0}\}$, an affine space with $V_0 = \{\vec{0}\}$ and $\dim(V_0) = 0$.

Assume the theorem holds at step $t$, with $\Omega_t = \alpha_t \oplus V_t$. We evaluate the action of an operation $O_t$:
\begin{itemize}[leftmargin=*]
    \item \textbf{Clifford Unitaries:} Completely absorbed by the rewinding map $\Phi_{t \to 0}$. They do not exist in the HIR and incur zero runtime cost. $\Omega_{t+1} = \Omega_t$, so $V_{t+1} = V_t$.
    \item \textbf{Non-Clifford Rotations:} A generic rotation decomposes via a Linear Combination of Unitaries into $c_1 I + c_2 P_{t \to 0}$. Applying this branches the superposition: $\Omega_{t+1} = \Omega_t \cup (\Omega_t \oplus \beta(P))$. The union of an affine space with its translation forms a new affine space.
    \begin{itemize}
        \item \textit{Collision:} If $\beta(P) \in V_t$, the rewound shift vector is already inside the active space. The space folds onto itself (an in-place interference butterfly). $\dim(V_{t+1}) = \dim(V_t)$.
        \item \textit{Expansion:} If $\beta(P) \notin V_t$, the underlying vector space gains exactly one basis vector. $\dim(V_{t+1}) = \dim(V_t) + 1$.
    \end{itemize}
    \item \textbf{Projective Pauli Measurements:} Measuring an observable applies a projector mapping to $I \pm M_{t\to0}$.
    \begin{itemize}
        \item \textit{Diagonal Filter ($\beta(M) = \vec{0}$):} The rewound observable acts diagonally. The projector enforces a linear parity constraint $\gamma(M) \cdot \alpha = c \pmod 2$. Intersecting the affine space $V_t$ with this hyperplane yields a new affine space. If it cuts the space, the dimension decreases by 1. Otherwise, it is unchanged.
        \item \textit{Internal Merge ($\beta(M) \in V_t \setminus \{\vec{0}\}$):} The measurement shifts the basis entirely within the active superposition, merging paired amplitudes. Factoring out the $\beta(M)$ direction yields the quotient space $V_{t+1} = V_t / \langle \beta(M) \rangle$. The dimension strictly decreases by 1.
        \item \textit{External AG Pivot ($\beta(M) \notin V_t$):} The measurement acts along a dimension entirely orthogonal to the active support. The topological frame update is handled purely by the deterministic AG row reductions captured in $\Phi_{t \to 0}$. Because the collapse occurs along a dimension unused by the active state ($\alpha_{\beta(M)} = 0 \;\forall \alpha \in \Omega_t$), the continuous amplitudes are mathematically isolated from the geometric collapse. The affine dimension is completely unaffected: $\dim(V_{t+1}) = \dim(V_t)$.
    \end{itemize}
\end{itemize}
By induction, $\Omega_t$ is always an affine space whose dimension is entirely dictated by the algebraic dependence of the HIR masks. The bound $|\mathbf{v}_t| \le 2^{\dim(V_t)}$ holds universally.
\end{proof}

Because the evolution of $V_t$ relies solely on the static masks $\beta(P)$ generated by the state-aware HIR, the compiler Back-End can deterministically track the active affine space and extract the $\max_t \dim(V_t)$ without ever running a single quantum shot. This allows the runtime VM to statically allocate its dense complex array exactly once, circumventing the tree-based dynamic memory overheads of traditional sparse simulators.

\subsection{The GHZ Corollary: Surpassing Spatial Bounds}

The power of framing the state in the rewound static space is profoundly evident when evaluating highly entangled quantum states, where static bounds inherently fail.

\begin{corollary}[The GHZ Collapse]
\label{cor:ghz}
Let $|\phi\rangle$ be an $n$-qubit GHZ state, and let $O_t$ be a sequence of $T$-gates applied to all $n$ qubits. Static spatial bounds predict $\mathcal{O}(2^n)$ complexity. The Dynamical Shift-Rank Bound proves $|\mathbf{v}| \le 2$, a constant bounded complexity entirely independent of $n$.
\end{corollary}

\textit{Proof:} Static bounds rely on the spatial footprint $|Q| = n$. Because no local $Z$ operator commutes with the global $X^{\otimes n}$ stabilizer of the GHZ state, the number of local $Z$-stabilizers $r_Q = 0$, falsely predicting a maximal $|\mathbf{v}| \le 2^n$ blowup.

However, Theorem \ref{thm:shift_rank} evaluates the temporal collision of the \textit{rewound} operators. Because the qubits are maximally entangled by prior Clifford operations, the Heisenberg rewinding map $\Phi_{t \to 0}$ pulls the spatial shift vectors for all $n$ local $Z$ operators into the exact same algebraic vector over $\mathbb{F}_2^n$; i.e., $\beta(Z_1) = \beta(Z_2) = \dots = \beta(Z_n) = \beta$.

The first $T$-gate expands the dimension by 1 ($\beta \notin V_0$). Every subsequent $T$-gate introduces the exact same shift mask ($\beta \in V_1$), triggering an $\mathcal{O}(1)$ \textit{Collision} (Lemma 2). The affine space dimension remains completely flat. Thus, $V_t = \text{span}(\beta)$, the dynamical rank is exactly 1, and $|\mathbf{v}| \le 2^1 = 2$.
% ==============================================================================
\section{The Physics of the Affine Space: The Active Magic Volume}
\label{sec:physics}
% ==============================================================================

The dimension $k(t) = \dim(V_t)$ is not merely a computational bound; it is a fundamental physical metric characterizing the non-classicality of the circuit, closely related to Stabilizer Nullity \cite{beverland2020} and the resource theory of magic states.

We define $k(t)$ as the \textbf{Active Magic Volume} of the circuit at time $t$. If $k(t)=0$, the system is strictly within the classical Gottesman-Knill regime. As $k$ grows, it measures the exact number of independent, non-classical orthogonal dimensions the state is actively exploring.

While raw $T$-count is a common proxy for simulation cost, it monotonically increases and vastly overestimates true complexity. The trajectory of $k(t)$ behaves as a discrete Markov chain—a "tug-of-war" of contextuality. Non-Clifford operations (like $T$-gates) attempt to inject contextuality, stepping $k$ up. Syndrome extractions and projective measurements actively cool and dissipate contextuality, stepping $k$ down.

A fault-tolerant quantum computing (FTQC) protocol may possess hundreds of $T$-gates, but if its peak rank $\max_t k(t)$ never exceeds a small integer, the circuit possesses very little true "quantumness." The magic injected by the non-Clifford gates is continually localized and collapsed by the error-correcting topology. The compiler's peak rank metric extracts the precise upper bound of a circuit's true quantumness, formalizing the boundary between classically simulable protocols and true quantum advantage.

% ==============================================================================
\section{The UCC Middle-End: Algebraic Optimization}
\label{sec:compiler}
% ==============================================================================

\textit{[Drafting Outline: Expand with specific programmatic details from UCC]}

\begin{itemize}[leftmargin=*]
    \item \textbf{The Heisenberg IR (HIR):} Explain how the AOT front-end absorbs Cliffords and emits a flat list of abstract Pauli operations defined by $\beta$ and $\gamma$ masks, formulating a global phase polynomial natively.
    \item \textbf{Symplectic Barriers for FTQC-Safe Compilation:} Detail how current $T$-count optimizers (e.g., PyZX, TODD) are ``FTQC-blind.'' They commute $T$-gates past physical noise channels to minimize gate counts, conjugating the physical errors (e.g., $TXT^\dagger \propto Y$) and destroying the hardware's exact error model required by QEC decoders. UCC evaluates the $\mathcal{O}(1)$ bitwise symplectic inner product between the rewound masks of an operation and stochastic nodes. This allows the optimizer to algebraically block illegal commutations, guaranteeing strictly \textit{Noise-Preserving} phase polynomial optimization.
    \item \textbf{Algebraic Elimination of the Synthesis Penalty:} Explain how mapping optimized phase polynomials back to CNOT ladders destroys hardware routing locality. UCC avoids this entirely by exporting the Optimized HIR directly to Pauli-Based Computation (PBC) hardware routers (e.g., lattice surgery), achieving zero synthesis overhead.
\end{itemize}

% ==============================================================================
\section{The UCC Virtual Machine: Systems Execution}
\label{sec:systems}
% ==============================================================================

\textit{[Drafting Outline: Expand with VM architecture specifics]}

\begin{itemize}[leftmargin=*]
    \item \textbf{Bounded Numerical Drift via Dominant Term Factoring:}
    Standard LCU simulation suffers from exponential IEEE-754 drift due to symmetric branching. We introduce a \textit{Bounded-Norm Gauge Choice}. By parameterizing non-Clifford gates as $T = \cos(\pi/8)[I - i\tan(\pi/8)Z]$ and factoring out the dominant Identity term (weight exactly $1.0$), classical physics is computed for free (0 FLOPs, 0 memory writes). Because relative spawned branches carry weights $\le \tan(\pi/8) \approx 0.414 < 1.0$, the $L_2$ variance is mathematically bounded, structurally preventing the exponential floating-point drift endemic to standard LCU simulators.
    \item \textbf{Cache-Aligned Instruction Stream:} Discuss the 32-byte cache-aligned bytecode opcodes (\texttt{OP\_BRANCH}, \texttt{OP\_COLLIDE}, \texttt{OP\_AG\_PIVOT}) and how $k_{\max}$ pre-allocates exactly one contiguous dense array per shot, eliminating dynamic memory overhead.
    \item \textbf{Geometric Gap Sampling:} Explain how Pauli noise is simulated without branching the instruction stream.
\end{itemize}

% ==============================================================================
\section{Demonstrations and Workflows}
\label{sec:demos}
% ==============================================================================

\textit{[Drafting Outline: Placeholder for benchmarking plots and results]}

\begin{itemize}[leftmargin=*]
    \item \textbf{Ground-Truth $d=5$ Magic State Cultivation:}
        \begin{itemize}
            \item Compare UCC's CPU-based performance to SOFT's 16-GPU H800 cluster benchmark. By taking advantage of the low peak dynamical rank, UCC sidesteps the massive GPU clusters required by non-dynamical simulators.
            \item Validate logical error rates against the empirical conjectures of Gidney et al. \cite{gidney2024magic}.
            \item Plot the Active Magic Volume $k(t)$ over the execution of the circuit to visually prove the rank-collapsing effects of the syndrome measurements.
        \end{itemize}
    \item \textbf{VQA Parameter Sweeps at Scale:} Show how mutating the LCU Constant Pool allows Variational Quantum Algorithms (QAOA/VQE) to update rotation angles with zero re-compilation latency.
    \item \textbf{Error Mitigation via Pauli Twirling:} Demonstrate that inserting random Pauli twirls has zero runtime overhead in UCC because they are absorbed into the branchless AOT Noise Schedule.
\end{itemize}

% ==============================================================================
\section{Conclusion}
\label{sec:conclusion}
% ==============================================================================

By mapping the generalized stabilizer formalism into an evolving affine subspace over $\mathbb{F}_2^n$, we have rigorously bound the complexity of universal quantum simulation. The Universal Compiler Collection (UCC) exploits Topological Determinism to separate heavy geometric tracking from probabilistic execution. This provides a unified framework capable of acting simultaneously as an FTQC-safe algebraic optimizer for hardware routing and a hyper-fast virtual machine for exact Monte Carlo simulation, providing an essential tool for the era of fault-tolerant quantum computing.

\begin{acknowledgments}
The authors acknowledge insightful discussions with [Names]. This research was supported by [Funding Agency].
\end{acknowledgments}

% ==============================================================================
\appendix
\section{The Active Magic Volume as a Physical and Computational Resource}
\label{app:amv}
% ==============================================================================

In Section \ref{sec:physics}, we defined the Active Magic Volume (AMV) at time $t$ as $k(t) = \dim(V_t)$, the dimension of the affine $\mathbb{F}_2$ subspace tracking the non-stabilizer degrees of freedom of the state. While traditional quantum complexity metrics rely on static operation counts (such as the total $T$-count or $T$-depth), the AMV provides an intensive, dynamical order parameter. In this appendix, we explore the theoretical implications of the AMV across fault-tolerant thermodynamics, structural complexity theory, and compiler optimization.

\subsection{The Thermodynamics of Fault Tolerance and Maxwell's Demon}

A central paradox in the classical simulation of fault-tolerant quantum computing (FTQC) is why highly parameterized subroutines—such as magic state distillation and cultivation, which contain dozens or hundreds of $T$-gates—remain classically tractable. Conversely, if FTQC actively suppresses all non-stabilizerness, how can a fault-tolerant quantum computer ever execute classically intractable algorithms to achieve quantum utility?

This paradox is resolved by decomposing the total Active Magic Volume into logical and transient components:
\begin{equation}
    k_{\text{total}}(t) = k_{\text{logical}}(t) + k_{\text{transient}}(t)
\end{equation}
Here, $k_{\text{logical}}$ tracks the algorithmic contextuality intentionally injected into the logical code block (the computational payload), while $k_{\text{transient}}$ tracks the non-stabilizerness generated by physical noise (e.g., coherent over-rotations) and the intermediate ``exhaust'' of distillation factories prior to purification.

In the framework of the Dynamical Shift-Rank Bound (Theorem \ref{thm:shift_rank}), a projective measurement $M$ only collapses a dimension of the magic space if its shift vector $\beta(M)$ lies strictly inside the active affine space $V_t$. When this occurs, the measurement structurally merges superposition branches, cooling the system ($k \to k-1$).

Fault-tolerant quantum computing acts as a thermodynamic sink—a literal Maxwell's Demon—for non-stabilizerness. Syndrome measurements are specifically designed to commute with the logical operators of the code. Because they commute, the shift vector of a syndrome measurement is mathematically orthogonal to the active logical subspace: $\beta(M_{\text{synd}}) = \vec{0}$ with respect to the logical frame. Consequently, the syndrome measurement acts purely as a \textit{Diagonal Filter} (as derived in the proof of Theorem \ref{thm:shift_rank}), leaving $k_{\text{logical}}$ completely undisturbed. The logical payload is insulated and allowed to grow.

Conversely, physical errors and intermediate distillation waste anti-commute with the syndromes. For these transient dimensions, $\beta(M_{\text{synd}}) \in V_t \setminus \{\vec{0}\}$. The syndrome measurement fiercely slices through this transient space, triggering an \textit{Internal Merge} and pumping the entropy out of the system into the classical measurement record.

This thermodynamic perspective rigorously explains the tractability of QEC component simulation. In a distance-5 Magic State Cultivation factory, 71 of the 72 physical $T$-gates generate $k_{\text{transient}}$, which is aggressively cooled by dense parity checks, keeping $k_{\text{total}}$ bounded by a small constant. However, when these factories feed purified states into a logical algorithm, $k_{\text{logical}}$ is allowed to grow monotonically, eventually surpassing the limits of classical tractability and achieving true quantum utility.

\subsection{Complexity Theory: Auditing Quantum Advantage}

When analyzing near-term algorithms—such as the Quantum Approximate Optimization Algorithm (QAOA) or Measurement-Based Quantum Computing (MBQC) schemes—researchers frequently cite high total $T$-counts or non-Clifford operation counts as evidence of classical intractability.

The AMV demonstrates that an infinite $T$-count does not guarantee quantum advantage. By decoupling the static \textit{resource cost} of a circuit (how many $T$-gates are applied) from its \textit{dynamical complexity} (how much contextuality is simultaneously sustained), we can define a new, rigorously verifiable island of classical simulability.

\begin{corollary}[Simulability of Bounded AMV]
Any quantum circuit family whose Peak Active Magic Volume is strictly bounded by $\max_t k(t) \le \mathcal{O}(\log n)$ is efficiently classically simulable in bounded-error probabilistic polynomial time (BPP), regardless of an arbitrarily large or infinite total $T$-count.
\end{corollary}

If an algorithm's topology induces massive structural shift-vector collisions (as seen in Corollary \ref{cor:ghz}) or features dense mid-circuit measurements that continuously collapse the active space, its true quantumness is strictly bottlenecked. The AMV therefore serves as a deterministic, polynomial-time classical screener. By tracking $k(t)$ over the execution of a proposed algorithm, we can rigorously audit whether the procedure accesses an exponentially large Hilbert space or is simply simulating classical dynamics in disguise.

\subsection{The Coherent Magic Capacity (CMC) Benchmark}
Currently, the industry benchmarks Quantum Processing Units (QPUs) using Quantum Volume \cite{cross2019} or Linear Cross-Entropy Benchmarking (XEB) \cite{arute2019}. These metrics rely on random unitary circuits heavily dominated by Clifford dynamics. They test entanglement routing but fail to isolate a QPU's ability to sustain coherent non-Clifford contextuality.

We propose a novel hardware benchmark: \textbf{The Coherent Magic Capacity (CMC)}.

In the Pauli spectrum, the AMV $k$ dictates the exact spreading of Pauli expectation values. Theoretical work has shown that this smearing, quantified by the 2-Stabilizer R\'enyi Entropy (SRE) $M_2(\rho) = -\log_2 ( 2^{-n} \sum_{P} \langle P \rangle^4 )$, is a rigorous magic monotone \cite{leone2022}. In our formalism, the theoretical SRE of a pure state scales exactly linearly with the AMV $k$. Furthermore, SRE can be experimentally measured using randomized Classical Shadows \cite{huang2020}.

The CMC protocol is executed as follows:
\begin{enumerate}
    \item \textbf{The Oracle:} Using the UCC compiler's Magic Register Allocation, we intentionally engineer a suite of test circuits with precise, theoretically escalating Active Magic Volumes: $k_{\text{ideal}}(t) \in \{5, 10, 15, \dots\}$.
    \item \textbf{The QPU Run:} The circuits are executed on a physical QPU. By appending a layer of randomized single-qubit Clifford twirls prior to measurement, we statistically extract the empirical SRE to estimate $k_{\text{exp}}(t)$.
    \item \textbf{Evaluating Leakage:} In a noiseless computer, $k_{\text{exp}} = k_{\text{ideal}}$. On real hardware, decoherence acts as a statistical mixture, strictly destroying magic. As errors accumulate, the hardware's empirical magic ``leaks'' faster than $T$-gates can inject it.
\end{enumerate}

The QPU's \textit{Peak Magic Capacity} is the maximum $k$ at which the hardware's experimentally measured contextuality statistically matches the UCC compiler's ideal prediction. This provides the first benchmark explicitly testing a QPU's ability to sustain the precise computational resource responsible for quantum advantage.


\subsection{Compiler Theory: Magic Register Allocation}

Current state-of-the-art quantum compilers optimize for a single metric: the minimization of the total $T$-count. However, when a quantum circuit is intended for classical simulation, or for execution on a hybrid quantum-classical architecture that relies on tensor-networks for real-time decoding, the total $T$-count is irrelevant. The true computational bottleneck is the Peak AMV, $\max_t k(t)$.

In classical compiler theory, register allocation algorithms minimize the number of \textit{simultaneously active} variables to fit them into limited CPU registers. Quantum compilers must adopt a similar strategy for non-stabilizerness, performing ``Magic Register Allocation.''

Because the dimension of the affine space $V_t$ is determined exclusively by the $\mathbb{F}_2$ commutativity of rewound Pauli shifts and measurements, a compiler can reorder operations to actively shape the $k(t)$ curve. In this paradigm, uncomputed magic behaves as a memory leak: every non-Clifford expansion allocates ``magic memory,'' and every merging measurement whose shift vector lies within the active space acts as garbage collection.

By strategically delaying $T$-gates until just before they are strictly required by a stabilizer measurement, or by eagerly pulling commuting measurements forward to collapse active branches as early as possible, a compiler can drastically suppress the Peak AMV. It is entirely possible for two logically equivalent circuits with identical $T$-counts to exhibit vastly different complexities: a naive scheduling may yield a Peak AMV of 50 (rendering it classically intractable), while an AMV-optimized schedule suppresses the peak to 12. Minimizing Peak AMV introduces a critical new optimization target for quantum compilers bridging the gap between quantum topologies and classical processing limits.
% ==============================================================================
% Bibliography
% ==============================================================================

\begin{thebibliography}{99}

\bibitem{gottesman1998} D. Gottesman, \textit{Theory of fault-tolerant quantum computation}, Phys. Rev. A \textbf{57}, 127 (1998).

\bibitem{bravyi2016} S. Bravyi and D. Gosset, \textit{Improved classical simulation of quantum circuits dominated by Clifford gates}, Phys. Rev. Lett. \textbf{116}, 250501 (2016).

\bibitem{yoder2012} T. J. Yoder, \textit{The generalized stabilizer formalism for simulating arbitrary quantum circuits} (2012), arXiv:1212.6253.

\bibitem{li2025soft} R. Li \textit{et al.}, \textit{SOFT: a high-performance simulator for universal fault-tolerant quantum circuits} (2025), arXiv:2512.23037.

\bibitem{aaronson2004} S. Aaronson and D. Gottesman, \textit{Improved simulation of stabilizer circuits}, Phys. Rev. A \textbf{70}, 052328 (2004).

\bibitem{beverland2020} M. Beverland \textit{et al.}, \textit{Lower bounds on the non-Clifford resources for quantum computations}, PRX Quantum \textbf{1}, 020321 (2020).

\bibitem{gidney2024magic} C. Gidney, N. Shutty, and C. Jones, \textit{Magic state cultivation: growing T states as cheap as CNOT gates} (2024), arXiv:2409.17595.

\end{thebibliography}

\end{document}
